% Physical unit 011; accepted Session 02.
\GordanUnitStart{011}{printed page 6}
\GordanTopFolio{6}
\[
 a\cos l\,dl\sqrt{1-e^2\cos^2l}=-\sqrt{\sin^2l_0-\sin^2l}\,ds,
\]
\[
 \cos l_0\,dl\sqrt{1-e^2\cos^2l}=\mp\sqrt{\sin^2l_0-\sin^2l}\,\cos l\,d\psi,
\]
quarum in posteriori signum aut negativum aut positivum ponendum est, prout longitudines, puncto $\Pi$ ab exordio recedente, aut augentur, aut minuuntur. Integrando fit:
\[
 s=-a\int_0^l {dl\cos l\sqrt{1-e^2\cos^2l}\over\sqrt{\sin^2l_0-\sin^2l}},
\]
\[
 \psi=\mp\int_0^l {dl\cos l_0\sqrt{1-e^2\cos^2l}\over\cos l\sqrt{\sin^2l_0-\sin^2l}}.
\]

\begin{center}\S\ 7.\end{center}
Definito limite $l$ ita, ut habeat valorem et majorem quam \gq{$-l_0$} et minorem quam \gq{$+l_0$} ex his formulis uniter apti valores et extensionis et longitudinis deducuntur. Cui consentaneum est punctum $\Pi$ in motione circulum $l$ uno tantum in loco secare, quod quia punctus constanti celeritate perpetuo movetur, fieri nequit. Qua de causa formulae a nobis propositae pertinent non ad totam extensionem lineae geodeticae $l_0$, sed ad eam partem, quae sita est inter exordium et punctum, quo tangitur circulus $-l_0$.

Ut igitur illae transformentur in formulas universi generis, pro latitudine $l$, quam hactenus per se variabilem posuimus, ejusmodi inducenda est variabilis, quae pro ratione puncti $\Pi$ ab exordio recedentis, crescat. Cujusmodi est angulus $\varphi$, qui definitur aequatione:
\[
 \sin l_0\cos\varphi=\sin l.
\]

\begin{center}\S\ 8.\end{center}
Recte angulus $\varphi$ continet eas, quas quaerimus, qualitates. Est enim ratio inter lineam geodeticam et angulum $\varphi$ haec, ut in exordio $\varphi=0$ sit, in sectione aequatoris $\varphi={\pi\over 2}$; in tactione circuli $-l_0$, $\varphi=\pi$, secunda in sectione aequatoris $\varphi={3\over2}\pi$, secunda in tactione circuli $l_0$, $\varphi=2\pi$ etc. Similiter linea procedens, cum tangit utrumque circulum in punctis iis, quae nominare placet puncta conversionis, est $\varphi=m\pi$; cum secat aequatorem, $\varphi={2m+1\over2}\pi$.

\begin{center}\S\ 9.\end{center}
Quoniam cuivis puncto lineae geodeticae $l_0$ certi quidam valores anguli $\varphi$ respondent, quantitates
\[
 l\qquad B\qquad T\qquad s\qquad \psi
\]
